\begin{table}[htbp]
\centering
\caption{Sample characteristics (N = 1,027)}
\label{tab:sample_char}
\begin{threeparttable}
\begin{tabular}{llr}
\toprule
\textbf{Variable} & \textbf{Category} & \textbf{N (\%)} \\
\midrule
\textbf{Age} & 18--34 & 217 (21.1\%) \\
             & 35--44 & 320 (31.2\%) \\
             & 45--54 & 231 (22.5\%) \\
             & 55--65 & 136 (13.2\%) \\
             & 65 and above & 123 (12.0\%) \\
\addlinespace
\textbf{Gender} & Female & 545 (53.1\%) \\
                & Male & 482 (46.9\%) \\
\addlinespace
\textbf{Education} & No formal education & 41 (4.0\%) \\
                   & Primary school & 128 (12.5\%) \\
                   & Middle school & 236 (23.0\%) \\
                   & High school & 217 (21.1\%) \\
                   & College or University & 287 (27.9\%) \\
                   & Graduate or above & 118 (11.5\%) \\
\addlinespace
\textbf{Urban--Rural} & Urban & 472 (46.0\%) \\
                      & Rural & 555 (54.0\%) \\

\bottomrule
\end{tabular}
\begin{tablenotes}[flushleft]
\footnotesize
\item Notes: Values are counts with row percentages in parentheses. Percentages may not sum to 100 due to rounding.

\end{tablenotes}
\end{threeparttable}
\end{table}
